\section{Higher odd degrees for the arithmetic quartic}
\label{d2ext:higher-odd-quartic}

The arithmetic root lattice also detects coefficients in every higher
odd degree. In this section write $A_2=A$ for the quartic ring of
Theorem~\ref{d2:thm:geometric}, let
$W_{\mathrm{ar}}=\mathbf Z_{(5)}[\rho]$ and
$E_{\mathrm{ar}}=\mathbf Q(\rho)$, and let
$i:V(x)\hookrightarrow\operatorname{Spec}A_2$ be its regular divisor.
We retain the root classes $\eta_g$ and their real Chern classes $e_g$
on the split affine surface $\mathcal Q$. Throughout this section
$g\in\mathbf F_5^\times$, and $\sigma_g$ sends $\rho$ to $\rho^g$.

\begin{theorem}
\label{d2ext:twodimalloddgeometric}
For the quartic arithmetic local ring $A_2$ and every $n\ge2$, the kernel of
$K_{2n-1}(A_2)\to K_{2n-1}(\Frac A_2)$ contains an actual infinite-order
class. The rational kernel has dimension at least two for even $n$
and at least one for odd $n$.
\end{theorem}
\begin{proof}
Borel's theorem gives $\dim_{\Q}K_{2n-1}(E_{\mathrm{ar}})_{\Q}=2$ and
$\dim_{\Q}K_{2n-1}(\Q)_{\Q}=0$ for even $n$, respectively $1$
for odd $n$ \cite[Proposition~12.2]{Borel}. Denote restriction and transfer by $\operatorname{res}$
and $\operatorname{Tr}$, respectively; they satisfy
$\operatorname{res}\operatorname{Tr}=\sum_g\sigma_g$ and
$\operatorname{Tr}\operatorname{res}=4$. Thus the invariant subspace
is exactly the restriction from $\Q$, and its trace-zero complement
$V_n$ has the asserted dimension. Quillen localization \cite{Quillen,Weibel} gives
$K_{2n-1}(W_{\mathrm{ar}})_{\Q}\simeq K_{2n-1}(E_{\mathrm{ar}})_{\Q}$.
Choose an actual $\gamma\in K_{2n-1}(W_{\mathrm{ar}})$ with noninvariant Borel
image and put $\beta_n=4\gamma-\sum_g\sigma_g\gamma$.
Its regulator is nonzero and its Galois sum is zero exactly.

Set $c_n=i_*\beta_n$. Flat base change again identifies its image
on $\mathcal T$ with $\sum_g(\sigma_g\beta_n|_{E_{\mathrm{ar}}})\eta_g$.
At a complex embedding, represent the point regulators
$b_g\in\C/\R(n)$ in the complementary real line $i\R(n)$
\cite[Section~6]{BunkeTamme}.
They form a nonzero vector with sum zero, hence a nonconstant vector.
For each finite open $U=\mathcal Q_s$, the product regulator is
represented by $2\pi i\sum_gb_ge_g$ in
\[
 \HD^3(U,\R(n+1))=H^2(U,\C)/H^2(U,\R(n+1)),
 \qquad n+1\ge3.
\]
Multiplication by $2\pi i$ induces an isomorphism
$\C/\R(n)\simeq\C/\R(n+1)$. The root lattice and
Proposition~\ref{d2:prop:geometry} therefore make this class
nonzero. If an integer killed $c_n$, continuity for the filtered
arithmetic localization would make that equality hold on some
$\mathcal Q_s$, contradicting the regulator. Localization at the
source divisor gives integral vanishing on $A_2[1/x]$.
Applying the argument to every nonzero vector of $V_n$ proves
injectivity on $V_n$ and gives the asserted rank bounds.
\end{proof}

The restriction $n\ge2$ is essential here: at $n=1$, the target
twist is two and $F^2$ on a complex affine surface is not zero.
The displayed quotient formula cannot be used in that degree.

\section{An explicit product of three units}
\label{d2ext:sec-explicit-triple}

Retain $A_2=A$ from Setup~\ref{d2:setup:ring} and its
completion $\widehat A_2=\widehat A$ from Theorem~\ref{d2:thm:completed}.
We identify an explicit product of three units with the weighted
root class after pullback to the split surface and the completed test algebra.

Let $N$ be a complex manifold of dimension two. Write $A^\bullet(N)$
for complex smooth forms and
\[
 I(p)^k=\{\xi\in A^k([0,1]\times N):
 \xi|_0\in (2\pi i)^pA^k_{\mathbf R}(N),\quad
 \xi|_1\in F^pA^k(N)\}.
\]
Thus $H^k(I(p))=H^k_{\mathcal D,\mathrm{an}}(N,\mathbf R(p))$,
with multiplication induced by wedge product.

\begin{lemma}\label{d2ext:interval-product}
Let $g$ be a nowhere-vanishing holomorphic function with a holomorphic
logarithm $L$, and let $z\in H^2(I(2))$. If $\xi$ is a closed
representative of $z$, put $\omega=\xi|_1$. Then $\omega$ is a
closed holomorphic two-form, is independent of the representative,
and
\[
 r_1(g)z=[-L\omega]\quad\hbox{in}\quad
 H^3_{\mathcal D,\mathrm{an}}(N,\mathbf R(3))
 =H^2(N,\mathbf C)/H^2(N,\mathbf R(3)).
\]
Here $r_1(g)$ is the analytic Deligne first regulator of $g$, with
the interval and cone conventions of
\cite[Definition~2.7 and Proposition~2.8]{BunkeTamme}.
\end{lemma}

\begin{proof}
Write $L=a+i\phi$, with $a,\phi$ real, and let $\tau$ denote the
interval coordinate. The usual representative of $r_1(g)$ is
\[
 i\,d\phi+d(\tau a).
\]
Its difference from $d(\tau L)$ is
$d(i(1-\tau)\phi)$. The primitive $i(1-\tau)\phi$ belongs to
$I(1)^0$: its value at $0$ is purely imaginary, hence in
$\mathbf R(1)$, and its value at $1$ is zero. Consequently
$d(\tau L)$ represents $r_1(g)$ in $I(1)$.

The endpoint $\omega=\xi|_1$ is closed and lies in $F^2A^2(N)$,
so it is a closed holomorphic two-form. Changing $\xi$ by an
$I(2)$-boundary cannot change this endpoint, because
$F^2A^1(N)=0$.

The product has representative
\[
 \upsilon=d(\tau L)\wedge\xi=d(\tau L\xi).
\]
Its endpoint at $0$ vanishes. Its endpoint at $1$ is
$dL\wedge\omega$, which vanishes by complex dimension two.
Thus the cone map sends its class to
$[-\int_0^1\upsilon]$. Fiberwise Stokes gives
\[
 -\int_0^1\upsilon
 =-L\omega+d_N\!\left(\int_0^1\tau L\xi\right),
\]
and hence gives the asserted representative. It is closed, since
$d(L\omega)=dL\wedge\omega=0$. The primitive $\tau L\xi$
has endpoint $L\omega$ at $1$; as $F^3A^2(N)=0$, it generally
does not belong to $I(3)^2$.

Finally $F^3A^\bullet(N)=0$, so the cone long exact sequence gives
the stated quotient: the next real-to-complex cohomology map is
injective. This proves the formula.
\end{proof}

\begin{proposition}
\label{d2ext:dimtwoExplicitTriple}
The triple
\begin{equation}\label{d2ext:eq-explicit-triple}
 c_{\rm exp}=\left\{2+y,\ 1-x,\
 1+(1-x)(x^3+Ty^3)y\bigl((1+y)^5+1\bigr)\right\}\in K_3(A_2)
\end{equation}
has infinite order in $A_2$ and in $\widehat A_2$, and its
restriction to the localization at $x$ is zero integrally in both cases.
\end{proposition}

\begin{proof}
\emph{Integral vanishing.}
Put $t=-1-y$, $v=x^3+Ty^3$, and
\[
 H(t)=\frac{t^{10}-1}{\Phi_{10}(t)}=(t^5-1)(t+1)
      =y\bigl((1+y)^5+1\bigr),\quad b=vH(t).
\]
Then $xv=-\Phi_{10}(t)$, so $1-xb=t^{10}$.
Set $g=1-t$, $\lambda=1-x$, and $\mu=1+\lambda b$.
These are units: their residues are $2,1,1$ respectively, and
$t$ is also a unit. In the convention
$\langle r,s\rangle=\{r,1-rs\}$ when $r$ is a unit, the
Dennis--Stein relations give
\[
 \delta=\langle x,b\rangle
 =\langle x,b+1-xb\rangle=\langle x,\mu\rangle
 =\{\lambda t^{10},\mu\}.
\]
Here $\langle x,1\rangle=0$ and $1-x\mu=\lambda t^{10}$;
see \cite[Chapter~III, Definition~5.11 and relations (D1)--(D3)]{Weibel}.
Since $\{g,t^{10}\}=10\{1-t,t\}=0$, multiplication gives
\begin{equation}\label{d2ext:dimtwoDSidentity}
 [g]\delta=\{g,\lambda,\mu\}=c_{\rm exp},\qquad
 [g]\delta|_{D(x)}=\{g,x,t^{10}\}=0.
\end{equation}
All identities are integral and persist under every ring map in use.

\emph{The universal supported coefficients.}
Write $E_{\mathrm{ar}}=\mathbf Q(\rho)$, where $\rho$ is primitive of order five,
and $t_g=-\rho^g$ for $g\in\mathbf F_5^\times$. On the smooth surface
\[
 Q^0=\operatorname{Spec}E_{\mathrm{ar}}[X,V,t]/(XV+\Phi_{10}(t)),\qquad
 U=Q^0[1/(t(1-t))],
\]
the class $\alpha=[1-t]\langle X,VH(t)\rangle$ is defined and
vanishes on $U[1/X]$. The entire divisor $D=V(X)$ is contained
in $U$. Localization and d\'evissage therefore give an actual
supported lift $\kappa\in K_3(Q^0)$ of $\alpha$.
Put $\eta_g=1-[(X,t-t_g)]$. Since the four components of $D$
are affine lines, $\kappa=\sum_g\gamma_g\eta_g$ with
$\gamma_g\in K_3(E_{\mathrm{ar}})$, and $\sum_g\eta_g=0$.

\emph{Independence of the supported lift.}
This lift is unique rationally as a class on $Q^0$.
Here is a direct $K$-theoretic check, independent of its periods,
using localization and the Laurent fundamental theorem \cite{Weibel}.
Put $Z=\mathbf A^1_{E_{\mathrm{ar}}}\setminus\{0,1\}$.  Then
$U[1/X]\simeq\mathbf G_m\times Z$, and localization on the
affine line, split by the units $t$ and $1-t$, gives
\[
 K_r(Z)\simeq K_r(E_{\mathrm{ar}})\oplus
       K_{r-1}(E_{\mathrm{ar}})[t]\oplus K_{r-1}(E_{\mathrm{ar}})[1-t].
\]
Together with the Laurent fundamental theorem this gives
\[
\begin{split}
 K_4(U[1/X])\simeq{}&
 K_4(E_{\mathrm{ar}})\oplus K_3(E_{\mathrm{ar}})[t]\oplus K_3(E_{\mathrm{ar}})[1-t]\\
 &{}\oplus K_3(E_{\mathrm{ar}})[X]
 \oplus K_2(E_{\mathrm{ar}})[X][t]\oplus K_2(E_{\mathrm{ar}})[X][1-t].
\end{split}
\]
The first three summands extend to $U$, so their boundary in
$K_3(D)=K_3(E_{\mathrm{ar}})^4$ is zero.  For a number field,
$K_2(E_{\mathrm{ar}})_{\mathbf Q}=K_4(E_{\mathrm{ar}})_{\mathbf Q}=0$: localization
reduces this to the rational even $K$-group vanishing for its
number ring \cite[Proposition~12.2]{Borel} and the torsion odd
$K$-groups of finite fields
\cite{QuillenFinite}\cite[IV, Corollary~1.13]{Weibel}.
The only remaining boundary is that of
$\beta[X]$, namely the diagonal tuple
$\pm(\beta,\beta,\beta,\beta)$ with a common convention-dependent
sign, since $X$ has order one on each root component. Localization therefore
identifies the kernel as
\[
 \ker\bigl(K_3(D)_{\mathbf Q}\longrightarrow K_3(U)_{\mathbf Q}\bigr)
       =\Delta K_3(E_{\mathrm{ar}})_{\mathbf Q}.
\]
The diagonal pushes forward to zero on $Q^0$, by the
length-one free resolution of the principal divisor $X=0$.
Consequently any two actual supported lifts of $\alpha$ give
the same class in $K_3(Q^0)_{\mathbf Q}$.
Moreover, since
$Q^0[1/X]\simeq\mathbf G_m\times\mathbf A^1_{E_{\mathrm{ar}}}$, the same
Laurent calculation gives
$\ker(K_3(D)\to K_3(Q^0))=\Delta K_3(E_{\mathrm{ar}})$ already integrally.
Thus its supported coefficients are unique modulo the diagonal.

Write $\operatorname{bl}$ for the Bloch-symbol map to algebraic
$K_3$ in \cite[Theorem~4.9]{BunkeTamme}. For a root $a$ of order ten define
$\beta_{\rm BT}(a)=\frac1{10}\operatorname{bl}(10[a])$
using \cite[Theorem~4.9 and (4.8)]{BunkeTamme}. Both $a$ and
$1-a$ are global units, and $10[a]$ has zero wedge boundary.
The normalization at every complex embedding $\sigma$ is
$r_2(\beta_{\rm BT}(a))_\sigma=-iD_{\rm BW}(\sigma a)$.
We claim the exact rational identity
\begin{equation}\label{d2ext:dimtwoUniversalCoefficients}
 \kappa=-10\sum_g\beta_{\rm BT}(t_g)\eta_g
          \quad\hbox{in }K_3(Q^0)_{\mathbf Q}.
\end{equation}

\emph{Determination of the supported coefficients.}
Fix a complex embedding $\sigma:E_{\mathrm{ar}}\hookrightarrow\mathbf C$ and
write $a_g=\sigma t_g$. The four roots are
$e^{\pm i\pi/5}$ and $e^{\pm3i\pi/5}$. For any distinct pair
$a_i,a_j$, their chord avoids $0$, $1$, and the other roots. Choose
one neighborhood of the chord on which the logarithm
$L=\operatorname{Log}(1-t)$ and the classical dilogarithm
$F=\operatorname{Li}_2(t)$ are single-valued. We use the branches
continued from their power series on the open unit disc; they satisfy
$dF=-L\,dt/t$.

Over this chord the matching sphere is parametrized by
\[
 X=\sqrt{|\Phi_{10}(t)|}\,e^{i\vartheta},\qquad
 V=-\frac{\Phi_{10}(t)}{|\Phi_{10}(t)|}\,\overline X,
\]
with $X=V=0$ at the two endpoints. This parametrization satisfies
$XV=-\Phi_{10}(t)$ and is smooth at each pole: $|X|^2$ is a
smooth positive multiple of the distance to that endpoint along
the chord. Orient the sphere by $ds\wedge d\vartheta$, where $s$
increases from $a_i$ to $a_j$.

On $X\ne0$, the weight-two holomorphic endpoint of the regulator
of $\delta=\langle X,VH(t)\rangle$ is
\[
 \omega=10\frac{dX}{X}\wedge\frac{dt}{t}
        =-\frac{dX\wedge d(VH(t))}{t^{10}}.
\]
The last expression extends holomorphically across both poles.
It therefore identifies the endpoint on the entire neighborhood.
Lemma~\ref{d2ext:interval-product} now gives
$r_3(\alpha)=[-L\omega]$. In particular this calculation retains
the full interval regulator, including its real endpoint.
On the punctured sphere,
$L\omega=10d(F\,dX/X)$. The boundary circles at the two poles
contribute $2\pi iF(a_j)$ and $-2\pi iF(a_i)$, and the integrals
over shrinking smooth caps tend to zero. The resulting period is
\begin{equation}\label{d2ext:matching-period}
 \int_{\Sigma_{ij}}r_3(\alpha)
   =-20\pi i\bigl(F(a_j)-F(a_i)\bigr)
       \pmod{\mathbf R(3)}.
\end{equation}
Since $|a_i|=|a_j|=1$, the real part is
$20\pi(D_{\rm BW}(a_j)-D_{\rm BW}(a_i))$.
The root divisor has local normal coordinate $X$. Its intersection
numbers with the oriented sphere are $+1$ at the initial pole
and $-1$ at the terminal pole. Thus the Chern periods of
$\eta_i$ and $\eta_j$ are $2\pi i$ and $-2\pi i$, respectively.
Multiplicativity gives, for every $\sigma$,
\[
 r_2(\gamma_i-\gamma_j)_\sigma
   =-10r_2\bigl(\beta_{\rm BT}(t_i)-\beta_{\rm BT}(t_j)\bigr)_\sigma.
\]
Borel injectivity \cite[Section~6]{BunkeTamme}, applied to all
embeddings, proves equality in $K_3(E_{\mathrm{ar}})_{\mathbf Q}$. The remaining coefficient is common to
all four components and contributes zero because $\sum_g\eta_g=0$.
This proves \eqref{d2ext:dimtwoUniversalCoefficients}.

\emph{Comparison with the actual cyclotomic coefficient.}
Write $\beta_{\rm BDJ}(a)$ for the field class supplied by the
symbol map of \cite[Theorem~1.10(1)--(2) and Remark~1.7]{BesserDeJeu}.
The complex formula of \cite[Proposition~4.1 and Remark~5.2]{DeJeuWedge}
compares its regulator with that of $\beta_{\rm BT}(a)$ by one
nonzero real scalar, independent of the field, root, and embedding.
This scalar is rational. Indeed, take $a=\zeta_6$ in
$L=\mathbf Q(\sqrt{-3})$ and put
$\beta_{\rm BT}(a)=\frac16\operatorname{bl}(6[a])$.
Both classes belong to the one-dimensional rational vector space
$K_3(L)_{\mathbf Q}$ \cite[Proposition~12.2]{Borel}
and are nonzero, since $D_{\rm BW}(\zeta_6)>0$. Consequently
$\beta_{\rm BDJ}(\zeta_6)=s\beta_{\rm BT}(\zeta_6)$ for a
unique $s\in\mathbf Q^\times$. Taking the complex regulator
identifies $s$ with the universal real scalar. Borel injectivity
at all embeddings \cite[Section~6]{BunkeTamme} then yields
\[
 \beta_{\rm BDJ}(a)=s\beta_{\rm BT}(a)
\]
for every root in use. The same field symbol map occurs in the
complex and $p$-adic formulas of the cited theorem.
Finite-residue localization identifies
$K_3(W_{\rm ar})_{\mathbf Q}$ with $K_3(E_{\mathrm{ar}})_{\mathbf Q}$,
so this comparison also holds before local completion.
Consequently the coefficient and character argument of
Proposition~\ref{d2:prop:coefficient} applies to
\[
 \Gamma=\operatorname{bl}(10[-\rho])\in K_3(\mathcal O_{E_{\mathrm{ar}}}),
 \qquad \beta=\Gamma-\sigma_{-1}\Gamma\in K_3(W_{\rm ar}).
\]
This is an actual integral antisymmetrization. Its residue class
vanishes because $\sigma_{-1}$ acts trivially on the residue field.
The identities $D_{\rm BW}(a^{-1})=-D_{\rm BW}(a)$ at every
embedding and Borel injectivity \cite[Section~6]{BunkeTamme} give
\[
 \beta_{\mathbf Q}=20\beta_{\rm BT}(-\rho)
                  =\frac{20}{s}\beta_{\rm BDJ}(-\rho).
\]
The syntomic symbol formula therefore makes the point regulator
of $\beta$ a common nonzero rational multiple of
$\operatorname{Li}_2(-\rho)$. The two congruences in
the proof of Proposition~\ref{d2:prop:coefficient} imply that its $e_1$ and
$e_3$ components are nonzero. The finite Chern factorization and
point comparison in Proposition~\ref{d2:prop:coefficient}
transfer this statement to its AMMN point coordinates.
Equation~\eqref{d2ext:dimtwoUniversalCoefficients} becomes
\begin{equation}\label{d2ext:dimtwoExactHalf}
 \kappa=-\frac12\sum_g\sigma_g(\beta)\eta_g.
\end{equation}

\emph{Detection on the arithmetic ring.}
On the chord $e^{-i\pi/5}\to e^{i\pi/5}$, the real period in \eqref{d2ext:matching-period} is
\[
 40\pi D_{\rm BW}(e^{i\pi/5})>0.
\]
The inequality follows from
$D_{\rm BW}(e^{i\theta})=-\int_0^\theta
\log(2\sin(v/2))\,dv$ for $0<\theta<\pi/3$.
This sphere lies in $Q^0[1/(t+1)]$, identified with $\mathcal Q$
by $X=x$, $V=x^3+Ty^3$, and $t=-1-y$.
The supported expression for $\kappa$ makes its regulator a real
root class. Proposition~\ref{d2:prop:geometry} keeps this class
nonzero on every permitted finite localization. If a nonzero
integer killed $c_{\rm exp}$, continuity would realize that
vanishing at one such finite localization, contradicting its
regulator. This proves infinite order in $A_2$ directly.

\emph{Detection after completion.}
The same substitutions define a ring map
$\mathcal O(Q^0)\to B[1/\pi]$,
with $X\mapsto\pi X$, $t\mapsto-1-\pi Y$.
Both $t$ and $1-t$ are units already in $B$, and the root line
$(X,t-t_g)$ pulls back to $J_g[1/\pi]$.
The denominator and completion maps of
Section~\ref{d2:sec:assembly} therefore give
\[
 c_{\rm exp}|_{B[1/\pi]}
 =\widehat c_{\rm exp}|_{B[1/\pi]}
 =-\frac12\sum_g\sigma_g(\beta)(1-[J_g])\big|_{B[1/\pi]}.
\]
Proposition~\ref{d2:prop:weighted} and the actual assembly
detect the right side as nonzero in rationalized completed
$K$-theory. Thus both actual source classes have infinite order.
The equality used here is on the split surface and its test pullback.
\end{proof}

\begin{corollary}
\label{d2ext:twodimexplicitMilnor}
For either $A=A_2$ or $A=\widehat A_2$, the ordinary Milnor symbol
\[
 m=\left\{2+y,\ 1-x,\
 1+(1-x)(x^3+Ty^3)y\bigl((1+y)^5+1\bigr)\right\}
       \in K_3^M(A)
\]
has infinite order and has zero image in
$K_3^M(\Frac A)$, integrally. In particular the ordinary
Milnor Gersten augmentation also fails after rationalization.
\end{corollary}

\begin{proof}
Write
\begin{align*}
 t&=-1-y,\qquad g=1-t,\qquad \lambda=1-x,\qquad q=t^{10},\\
 b&=(x^3+Ty^3)y\bigl((1+y)^5+1\bigr),\qquad \mu=1+\lambda b.
\end{align*}
The defining equation gives $1-xb=q$ and hence
$1-x\mu=\lambda q$. The elements $g,\lambda,\mu,q$
are units in $A$: their residues are respectively $2,1,1,1$.
In its fraction field, $x$ and $x\mu$ are also nonzero.
The two Steinberg relations and skew-symmetry in field
Milnor $K$-theory give
\begin{align*}
 0&=\{x\mu,\lambda q\},\qquad \{x,\lambda\}=0,\\
 \{x,q\}&=-\{\mu,\lambda\}-\{\mu,q\}
            =\{\lambda q,\mu\}.
\end{align*}
Furthermore $\{g,q\}=10\{1-t,t\}=0$. Therefore, without
rationalizing,
\[
 m=\{g,\lambda,\mu\}
   =\{g,\lambda q,\mu\}
   =\{g,x,q\}
   =-\{g,q,x\}=0
       \quad\text{in }K_3^M(\Frac A).
\]
The natural homomorphism $K_3^M(A)\to K_3(A)$ induced
by products of units sends $m$ to the already detected
infinite-order class $c_{\rm exp}$. Thus no nonzero integer
can kill $m$. 
\end{proof}

\section{A tame family at every residue prime}
\label{d2ext:sec-tame-family}

\begin{theorem}
\label{d2ext:twodimtamefamily}
Let $p$ be any prime and let $d\ge3$ satisfy $p\nmid d$. Put
\[
 A_{p,d}=\left(\Z_{(p)}[T,x,y]/
       (p+x^d+y^d+Txy^{d-1})\right)_{(p,x,y)}.
\]
This is a regular local ring of dimension two, with residue field
$\F_p(T)$, and
\[
 \dim_{\Q}\ker\bigl(K_3(A_{p,d})_{\Q}
       \longrightarrow K_3(\Frac A_{p,d})_{\Q}\bigr)
       \ \ge\ \lfloor d/2\rfloor.
\]
In particular these dimensions are unbounded with the residue
prime and local dimension fixed. The theorem concerns the
uncompleted rings displayed above.
\end{theorem}

\begin{proof}
\emph{The actual source and the flat splitting ring.}
Write $A=A_{p,d}$. Its regular ambient local ring has parameters
$p,x,y$, and its defining equation has linear term $p$.
Thus $A$ is regular of dimension two with maximal ideal $(x,y)$.
Solving the equation for $T$ gives $\Frac A=\Q(x,y)$.

Choose $\theta^d=-p$ and set
\[
 K=\Q(\theta),\qquad W_0=\Z_{(p)}[\theta],\qquad
 A/(x)=W_0[T]_{(\theta)}.
\]
Eisenstein's criterion gives $[K:\Q]=d$; $W_0$ is a DVR with
uniformizer $\theta$ and residue field $\F_p$.
Quillen localization gives
$K_3(W_0)_{\Q}\simeq K_3(K)_{\Q}$, and every actual
$K_3(K)$ class lifts to $K_3(W_0)$, since $K_2(\F_p)=0$.
Pull back such a coefficient $\beta$ to $A/(x)$ and push it
forward to obtain $c_\beta\in K_3(A)$. Its restriction to
$A[1/x]$ is zero. We will prove that the resulting rational map
$K_3(K)_{\Q}\to K_3(A)_{\Q}$ is injective.

Let $E=K(\mu_d)$ and choose a prime $\mathfrak P$ above $p$.
Put $W=(\mathcal O_E)_{\mathfrak P}$ and let $k$ be its finite
residue field. The completion of $E$ at this prime is
$\Q_p(\mu_d,\theta)$. Since $p\nmid d$, adjoining $\mu_d$
is unramified, and $Z^d+p$ remains Eisenstein over that
unramified extension. Consequently $\theta$ is a uniformizer
of $W$, and the $d$ roots of unity $(\zeta_i)_{i=1}^d$ in $W$ have distinct
reductions in $k$.
The ring $W$ is flat over $\Z_{(p)}$, but need not be finite.
Use the explicit local ring
\[
 B=(A\otimes_{\Z_{(p)}}W)_{(\theta,x,y)}
  =\left(V[x,y]/
      (x^d+y^d+Txy^{d-1}-\theta^d)\right)_{(\theta,x,y)},
 \qquad V=W[T]_{(\theta)}.
\]
The ideal indicated is maximal with residue field $k(T)$, and
$A\to B$ is flat by tensor product followed by localization.

\emph{The exceptional curve and its root points.}
The projectivized tangent cone of $B$ is
\[
 C:\ X^d+Y^d+TXY^{d-1}=U^d
       \subset\mathbf P^2_{k(T)}.
\]
It is geometrically smooth. A singular point must have $U=0$;
then $Y\ne0$ and $q(s)=s^d+Ts+1$ would have a multiple root.
The identity $dq-sq'=(d-1)Ts+d$ shows this can occur only at
finitely many algebraic values of $T$, or never if $p\mid d-1$.
Thus it cannot occur over $k(T)$.

The smooth plane curve $C$ is geometrically integral, so the
associated graded ring of $B$ is a domain. Krull intersection
then shows that $B$ is a domain. The blow-up
$\mathcal Y=\operatorname{Bl}_{(\theta,x,y)}\Spec B$
is regular and has the single reduced exceptional curve $C$.
Here are the local details needed for this assertion. Its
$\theta$-chart has equation
$X^d+Y^d+TXY^{d-1}=1$, with $x=\theta X$, $y=\theta Y$;
the other two strict-transform charts reduce to the corresponding
charts of the smooth $C$. A nonzero differential tangent to
the ambient exceptional plane therefore proves regularity at
every point of $C$. The nonregular locus is closed by excellence.
Its proper image in the local scheme $\Spec B$, if nonempty,
would contain the closed point, a contradiction. This proves
regularity everywhere.
Off the closed point the blow-up is an isomorphism; the
hypersurface $S_2$ property then also shows $B$ is normal.

The tautological ideal gives
$\mathcal O_{\mathcal Y}(C)|_C=\mathcal O_C(-1)$.
The root primes $(x,y-\theta\zeta_i)$ have strict transforms
meeting $C$ at $P_i=[1:0:\bar\zeta_i]$, each with multiplicity
one. Indeed, on the $\theta$-chart the $Y$-derivative is a unit
at $X=0,Y=\zeta_i$, and the two regular parameters there are
$\theta,X$.

\emph{The exceptional Picard lattice.}
On $C^*=C\cap\{UY\ne0\}$ the root-point classes have only
the diagonal relation. To prove this, extend $k$ to $\bar k$ and
consider
\[
 M=\Spec\bar k[X,Z,Y^{\pm1}]/(XZ-(1-Y^d)),\qquad
 T=(Z-X^{d-1})/Y^{d-1}.
\]
On inverting $X$ its ring is the Laurent UFD
$\bar k[X^{\pm1},Y^{\pm1}]$. Divisor localization and its units
$cX^aY^b$ give
$\Pic(M)=\Z^d/\Z(1,\ldots,1)$, generated by
$(X,Y-\bar\zeta_i)$.

Every finite fiber $T=t$ is geometrically integral. Under
$s=X/Y$, $w=1/Y$ it is
$w^d=s^d+ts+1$, $w\ne0$.
The polynomial $q_t=s^d+ts+1$ has $\deg\gcd(q_t,q_t')\le1$
by $dq_t-sq_t'=(d-1)ts+d$ and $d\ne0$ in $\bar k$.
If it were an $\ell$th power for a prime $\ell\mid d$,
that degree would be at least $d-d/\ell\ge2$, a contradiction.
A rational-function power would already be a polynomial power.
Since $\bar k$ contains $\mu_d$ and $p\nmid d$, this also
proves irreducibility of $w^d-q_t$: the root extension is
cyclic, and a proper degree would express $q_t$ as such a power.
Each removed $T$-fiber is therefore one principal prime divisor,
of multiplicity one. Divisor localization gives
$\Pic(M)\simeq\Pic(M_{\bar k(T)})$.
The generic fiber of $M$ is $C^*_{\bar k(T)}$, proving the root-point
assertion.

\emph{Arithmetic Picard and Betti persistence.}
The generic affine surface and test localization are
\begin{align*}
 Q&=\Spec E[T,x,y,y^{-1}]/(p+x^d+y^d+Txy^{d-1})\\
  &\simeq\Spec E[x,v,y^{\pm1}]/(xv+y^d+p),\qquad
 \mathcal T=\Spec B[1/\theta,1/y],
\end{align*}
where $v=x^{d-1}+Ty^{d-1}$. Inverting $x$ again gives a
Laurent UFD. Thus $\Pic(Q)=\Z^d/\Z(1,\ldots,1)$, generated
by $D_i=(x,y-\theta\zeta_i)$; this remains true after every
coefficient field extension.
The map $\Pic(Q)\to\Pic(\mathcal T)$ is injective.
Indeed, a root combination made principal on $\mathcal T$
can acquire on $\mathcal Y$ only strict divisors over $\theta=0$
or $y=0$, and a multiple of $C$. The former miss $C^*$ and
$\mathcal O_C(-1)$ is trivial on $U\ne0$.
Restriction to $C^*$ gives the same combination of the $P_i$,
so all coefficients are equal, which was already the zero
relation on $Q$.

Write $\mathcal O(\mathcal T)=\varinjlim_s\mathcal O(Q_s)$,
where $s$ runs over the permitted arithmetic denominators.
Each $E$-prime divisor removed by such an $s$ has zero Picard
class by the injection just proved. Its geometric components
are one Galois orbit. The absolute Galois group over $E$ acts trivially on the torsion-free
root Picard lattice, so each geometric component also has zero
class. The degree-two support sequence consequently gives
\[
 H^2(Q(\C),\R)\lhook\joinrel\longrightarrow
 H^2(Q_s(\C),\R)\qquad\text{for every permitted }s.
\]
The root Chern classes $e_i$ span an injected copy of
$\R^d/\R(1,\ldots,1)$: in the Gysin sequence of the $D_i$,
the angular class of $x$ has residue the diagonal and that of
$y$ has residue zero. This proves precisely the needed real
root-class persistence.

\emph{Detection of every nonzero coefficient.}
Let $\iota_i:K\hookrightarrow E$ send $\theta$ to
$\theta\zeta_i$, and put $\eta_i=1-[\mathcal O_Q(-D_i)]$.
Flat base change for the actual Gysin class and the projection
formula, followed by restriction to the regular $\mathcal T$,
give
\[
 c_\beta|_{\mathcal T}
    =\left.\sum_i\iota_i(\beta)\eta_i\right|_{\mathcal T}.
\]
At one complex embedding of $E$, the point regulators
$b_i\in\C/\R(2)$ run through all embeddings of $K$.
Choose their representatives in $i\R$. If $\beta\ne0$ in
$K_3(K)_{\Q}$, Borel injectivity \cite[Section~6]{BunkeTamme}
makes this vector nonzero.
Conjugate embeddings have opposite coefficients and real
embeddings have coefficient zero, so it is not diagonal.

The absolute multiplicative regulator of Bunke--Tamme
\cite[Theorem~2.31]{BunkeTamme} of the supported sum
is the real root class $2\pi i\sum_i b_i e_i$, modulo
$H^2(Q_s(\C),\R(3))$. It is nonzero on every permitted finite
open by the preceding injection. Here the analytic Deligne
target is $H^2(Q_s,\C)/H^2(Q_s,\R(3))$ because its complex
dimension is two and $F^3=0$; a nonzero real class survives this
quotient. Each individual regulator calculation spreads to a
regular finite type arithmetic model after inverting one
integer. The standard continuity
$K_3(\mathcal T)=\varinjlim_s K_3(Q_s)$ now rules out any
nonzero rational coefficient in the kernel of the Gysin map.

Finally $Z^d+p$ has one real root when $d$ is odd and none
when $d$ is even. Borel's rank calculation
\cite[Proposition~12.2]{Borel} therefore gives
$\dim_{\Q}K_3(K)_{\Q}=r_2(K)=\lfloor d/2\rfloor$.
Choosing actual lifts of a rational basis yields the asserted
number of independent classes of infinite order in the generic kernel.
All uses of localization, base change and regulators are the
standard inputs cited above \cite{Quillen,QuillenFinite,Borel,BunkeTamme,Weibel}.
\end{proof}

For example, $(p,d)=(2,3)$ gives
$\bigl(\Z_{(2)}[T,x,y]/(2+x^3+y^3+Txy^2)\bigr)_{(2,x,y)}$.
For any fixed $p$, arbitrarily large odd $d$ prime to $p$ give
the unbounded rank assertion without enlarging the residue
field $\F_p(T)$. 

\begin{remark}
The condition $d\ge3$ is used in the vertical-fiber calculation.
When $d=2$ and $p\ne2$, the fiber at $t=2$ has
$s^2+2s+1=(s+1)^2$, so the Kummer irreducibility argument is
unavailable and the vertical fiber splits. Localization can then
kill a root divisor difference. The proof therefore supplies no
degree-two analogue of the displayed tame family.
\end{remark}

\begin{corollary}
\label{d2ext:twodimtameallodd}
For $A_{p,d}$ as in Theorem~\ref{d2ext:twodimtamefamily} and every
$n\ge2$, the rational kernel of
$K_{2n-1}(A_{p,d})\to K_{2n-1}(\Frac A_{p,d})$ has dimension at least
\[
 \begin{cases}
 \lfloor d/2\rfloor,&n\text{ even},\\
 \lceil d/2\rceil-1,&n\text{ odd}.
 \end{cases}
\]
These bounds are realized by actual classes of infinite order.
For odd $d$ the bound is $(d-1)/2$ in every odd $K$-degree at least three.
\end{corollary}

\begin{proof}
Retain $K=\Q(\theta)$, $W_0$, the splitting field $E$, and the
$d$ embeddings $\iota_i:K\hookrightarrow E$ from the theorem.
Use the specified rational coefficient space
\[
 V_n=\ker\left(\operatorname{Tr}_{K/\Q}:
       K_{2n-1}(K)_{\Q}\longrightarrow K_{2n-1}(\Q)_{\Q}\right).
\]
For an actual $\gamma\in K_{2n-1}(K)$ set
$\beta=d\gamma-\operatorname{res}_{K/\Q}
                 \operatorname{Tr}_{K/\Q}\gamma$.
The identity $\operatorname{Tr}\operatorname{res}=d$ shows
that $\operatorname{Tr}\beta=0$ integrally. These actual
classes span $V_n$ rationally, because the displayed operator
acts as multiplication by $d$ on $V_n$.
Each $\beta$ lifts to $K_{2n-1}(W_0)$, since
$K_{2n-2}(\F_p)=0$. Localization, flat base change and the
projection formula give an actual class vanishing generically whose
test image is $\sum_i\iota_i(\beta)\eta_i$.

At a fixed complex embedding of $E$ let $b_i$ be the point
regulators, represented in the complementary real line
$i\R(n)$ to $\R(n)$ in $\C$. Base change for field transfer
gives $\sum_i\iota_i(\beta)=
\operatorname{res}_{E/\Q}\operatorname{Tr}_{K/\Q}\beta=0$;
therefore $\sum_i b_i=0$. If $\beta$ is nonzero rationally,
Borel injectivity \cite[Section~6]{BunkeTamme} makes the vector
$(b_i)$ nonzero and hence nondiagonal. Its line-bundle product is
$2\pi i\sum_i b_i e_i$. The preceding Betti persistence
keeps this root class nonzero on every permitted finite open.
The relevant analytic Deligne group is
\[
 H^3_{\mathcal D,\mathrm{an}}(Q_s,\R(n+1))
   =H^2(Q_s,\C)/H^2(Q_s,\R(n+1)),
\]
because $n+1\ge3$ and $Q_s(\C)$ is a surface.
Multiplication by $2\pi i$ identifies
$\C/\R(n)$ with $\C/\R(n+1)$, so the product regulator
is nonzero. The same continuity argument proves injectivity
of the supported map on $V_n$; applying it to every rational
relation gives independent actual representatives.

Rational transfer is surjective, since its composition with
restriction is multiplication by $d$. Borel's ranks
\cite[Proposition~12.2]{Borel} give
\[
 \dim V_n=
 \begin{cases}
 r_2(K),&n\text{ even},\\
 r_1(K)+r_2(K)-1,&n\text{ odd}.
 \end{cases}
\]
Here $r_2(K)=\lfloor d/2\rfloor$, while $r_1(K)$ is one for odd $d$ and zero for even $d$; the bounds follow.
\end{proof}

\section{A fixed arithmetic ring with infinite rank}
\label{d2ext:sec-strict-coefficients}

\begin{theorem}
\label{d2ext:twodimstrictcoefficientinfinite}
Fix a prime $p$ and an integer $d\ge3$ prime to $p$. Let
$V=\Z_{(p)}^{\mathrm{sh}}$, with residue field $\overline{\F}_p$,
and put
\[
 R=\left(V[T,x,y]/(p+x^d+y^d+Txy^{d-1})\right)_{(p,x,y)}.
\]
Then $R$ is a Noetherian regular local ring of dimension two,
with residue field $\overline{\F}_p(T)$. For every $n\ge2$,
\[
 \dim_{\Q}\ker\!\left(K_{2n-1}(R)_{\Q}
       \longrightarrow K_{2n-1}(\Frac R)_{\Q}\right)=\infty.
\]
The coefficient DVR $V$ is strictly henselian; no henselianity
of the displayed $R$ is asserted. This theorem concerns this
single uncompleted ring.
\end{theorem}

\begin{proof}
\emph{The coefficient and splitting rings.}
The strict henselization $V$ is a DVR with uniformizer $p$,
and $F=\Frac V$ is algebraic over $\Q$; see the Stacks Project,
\href{https://stacks.math.columbia.edu/tag/0AP3}{Tag 0AP3}
and the proof of
\href{https://stacks.math.columbia.edu/tag/09ET}{Tag 09ET}.
It is excellent because it is a DVR of fraction characteristic
zero
(\href{https://stacks.math.columbia.edu/tag/07QW}{Tag 07QW(4)}).
The simple roots of $Z^N-1$ in $\overline{\F}_p$ lift by
Hensel's lemma whenever $p\nmid N$, so $\mu_N\subset V$.
The ambient local ring defining $R$ has parameters $p,x,y$,
and the defining equation has initial term $p$. Thus $R$ has
the stated regularity, dimension and residue field.

Choose $\theta^d=-p$ and set
\[
 W=V[\theta],\qquad E_\infty=F(\theta),\qquad
 G=\operatorname{Gal}(E_\infty/F)=\{\sigma_\zeta:\zeta\in\mu_d\},
 \quad \sigma_\zeta(\theta)=\zeta\theta .
\]
Eisenstein shows that $W$ is a finite free DVR extension of
$V$, with uniformizer $\theta$, residue field $\overline{\F}_p$,
and degree $d$. Its displayed Galois action preserves $W$.
Moreover
\[
 D=\Spec(R/(x))=\Spec W[T]_{(\theta)}
\]
is regular. Pullback to $D$ followed by its Gysin map
$i_*:K_{2n-1}(D)\to K_{2n-1}(R)$ therefore constructs
actual classes vanishing on $R[1/x]$.

\emph{The geometric detector over $E_\infty$.}
Use the flat splitting map to
\[
 B=(R\otimes_VW)_{(\theta,x,y)},\qquad
 \mathcal T=\Spec B[1/\theta,1/y],
\]
and the smooth surface
\[
 Q=\Spec E_\infty[x,v,y^{\pm1}]/(xv+y^d+p),
 \qquad v=x^{d-1}+Ty^{d-1}.
\]
The geometric proof of Theorem~\ref{d2ext:twodimtamefamily}
applies over these coefficient rings. The required geometric calculations are as follows.
The blow-up still has the single smooth exceptional curve
\[
 C:\ X^d+Y^d+TXY^{d-1}=U^d
       \subset\mathbf P^2_{\overline{\F}_p(T)}.
\]
The same three charts prove its regularity along $C$, and
excellence and properness prove regularity everywhere. The
root strict transforms meet it transversely at
$[1:0:\bar\zeta]$. On $C\cap\{UY\ne0\}$ their Picard
classes have the lattice $\Z^d/\Z(1,\ldots,1)$:
the auxiliary surface
$XZ=1-Y^d$, $Y\ne0$, has that lattice, and every finite
$T=(Z-X^{d-1})/Y^{d-1}$ fiber is integral by the
$dq-sq'$ argument. These calculations already take place
over $\overline{\F}_p$ and do not require finite residue
constants.

Consequently, restriction of divisors to this exceptional open
gives
\[
 \Pic(Q)=\Z^d/\Z(1,\ldots,1)
       \lhook\joinrel\longrightarrow\Pic(\mathcal T),
\]
with generators $D_\zeta=(x,y-\zeta\theta)$.
The Picard calculation is unchanged by coefficient field
extension. Each $E_\infty$-prime divisor removed by a permitted
finite localization $Q_s$ has zero Picard class. Its geometric
components form one Galois orbit and all have the same class,
since the absolute Galois group over $E_\infty$ fixes its root
generators and the lattice is
torsion-free. Each component class is therefore zero.
The support sequence gives, for every embedding
$\jmath:E_\infty\hookrightarrow\C$,
\begin{equation}
 H^2(Q_\jmath(\C),\R)
      \lhook\joinrel\longrightarrow H^2((Q_s)_\jmath(\C),\R).
 \label{d2ext:strictcoefficientbetti}
\end{equation}
The Gysin sequence with complement $(\mathbf G_m)^2$ identifies
the root subspace with $\R^d/\R(1,\ldots,1)$, since the
angular classes of $x$ and $y$ have residues respectively the
diagonal vector and zero. This argument uses no regulator on
an infinite arithmetic base.

\emph{Finite actual coefficient spaces of unbounded rank.}
For $N>2$ divisible by $d$ and prime to $p$, choose compatible
roots of unity in $V$ and set
\[
 L_N=\Q(\mu_N)\subset F,\qquad K_N=L_N(\theta)\subset E_\infty .
\]
The chosen $p$-adic place of $L_N$ is unramified over $\Q_p$.
Eisenstein there gives $[K_N:L_N]=d$, and
$\operatorname{Gal}(K_N/L_N)$ is the restriction of $G$. Both fields
are totally imaginary. For fixed $n\ge2$, Borel's ranks
\cite[Proposition~12.2]{Borel} and restriction--transfer identities give
\begin{equation}
 V_{N,n}=
 \ker\!\left(\sum_{\sigma\in G}\sigma:
       K_{2n-1}(K_N)_{\Q}\to K_{2n-1}(K_N)_{\Q}\right),
 \qquad
 \dim_{\Q}V_{N,n}=\frac{(d-1)\varphi(N)}2.
 \label{d2ext:strictcoefficientrank}
\end{equation}
Indeed, the invariant subspace is exactly the restriction of
$K_{2n-1}(L_N)_{\Q}$, because
$\operatorname{res}\operatorname{Tr}=\sum_\sigma\sigma$ and
$\operatorname{Tr}\operatorname{res}=d$.
The two Borel ranks are $d\varphi(N)/2$ and $\varphi(N)/2$
in every degree under consideration.

Let $W_N=(\mathcal O_{K_N})_{\mathfrak P_N}$ be the DVR
at the prime induced by $W$. It maps to $W$, and $G$ preserves
it: above the chosen $L_N$ place the extension is totally
ramified and has a unique prime. Finite-field localization
gives
\[
 K_{2n-1}(W_N)_{\Q}\simeq K_{2n-1}(K_N)_{\Q},
\]
as well as actual surjectivity to the field, since
$K_{2n-2}$ of the finite residue field is zero.
Thus $V_{N,n}$ defines a rational linear map into
$K_{2n-1}(R)_{\Q}$ by coefficient pullback and $i_*$.
Every element of $V_{N,n}$ has an actual representative after clearing
denominators. If exact integral trace zero is wanted, replace
an actual lift $\gamma$ by
$d\gamma-\sum_{\sigma\in G}\sigma\gamma$; this operation
acts as multiplication by $d$ on $V_{N,n}$.

\emph{Detection, with both continuity passages explicit.}
Suppose $0\ne\beta\in V_{N,n}$ mapped to zero in
$K_{2n-1}(R)_{\Q}$. Borel injectivity \cite[Section~6]{BunkeTamme} supplies an embedding
$j_0:K_N\hookrightarrow\C$ at which its point regulator
is nonzero. Extend $j_0$ to an embedding
$\jmath:E_\infty\hookrightarrow\C$, using that
$E_\infty/K_N$ is algebraic.
For $\sigma\in G$, let $b_\sigma$ denote the point regulator
at $j_0\sigma$, represented in the complementary real line
$i\R(n)$ to $\R(n)$ in $\C$. These coefficients satisfy
\[
 b_1\ne0,\qquad \sum_{\sigma\in G}b_\sigma=0,
\]
so this vector is not diagonal. The detecting embedding is
chosen for this coefficient; no single embedding is required
to detect every element of $V_{N,n}$.

Put $\eta_\zeta=1-[\mathcal O_Q(-D_\zeta)]$. Flat Gysin
base change and the projection formula identify the image of
the supported class on $\mathcal T$ with the restriction of
\[
 \alpha=\sum_{\zeta\in\mu_d}(\sigma_\zeta\beta)\eta_\zeta
       \in K_{2n-1}(Q)_{\Q}.
\]
The assumed vanishing therefore makes $\alpha$ zero on
$\mathcal T$. First, $K$-theory continuity for the filtered
localization
$\mathcal O(\mathcal T)=\varinjlim_s\mathcal O(Q_s)$
gives a finite permitted denominator set $s$ for which
$\alpha|_{Q_s}=0$ rationally.
Second, $E_\infty$ is a filtered union of number fields.
All coefficients of these finitely many denominators, the
root ideals, and the chosen $K$-classes descend to a number
field $E_{\mathrm{ar}}\subset E_\infty$ containing $K_N$.
Continuity for this filtered coefficient extension implies
that the asserted vanishing already holds on the corresponding
finite open over some larger number field
$E_1\subset E_\infty$.

Apply the published absolute multiplicative regulator to that
finite stage, after spreading it to a regular finite-type
arithmetic model by inverting one integer. Evaluate at
$\jmath|_{E_1}$. The complex open of this finite model is exactly
$(Q_s)_\jmath$, and naturality gives the product class
\[
 2\pi i\sum_{\zeta\in\mu_d}b_{\sigma_\zeta}e_\zeta
 \quad\text{in}\quad
 H^3_{\mathcal D,\mathrm{an}}((Q_s)_\jmath,\R(n+1))
 =\frac{H^2((Q_s)_\jmath,\C)}
        {H^2((Q_s)_\jmath,\R(n+1))}.
\]
The equality holds because $n+1\ge3$ and the complex
dimension is two. Multiplication by $2\pi i$ identifies
$\C/\R(n)$ with $\C/\R(n+1)$. The non-diagonal root
vector and \eqref{d2ext:strictcoefficientbetti} make this
regulator nonzero, contradicting the finite-stage vanishing.
This proves injectivity on every $V_{N,n}$.
No commutation of analytic cohomology with an infinite
coefficient colimit has been used.

Finally, take for example $N=d^a$ as $a\to\infty$.
The dimensions in \eqref{d2ext:strictcoefficientrank} are
unbounded, while $R$ remains fixed. This proves infinite
rational rank in every degree $2n-1\ge3$.
All coefficient lifts, transfers and products used above are
actual Quillen $K$-theory operations
\cite{Quillen,QuillenFinite,Weibel,Borel,BunkeTamme}.
\end{proof}

For each $p$ one may fix $d=3$ if $p\ne3$, and $d=4$
if $p=3$. This gives one ring at each residue prime
with the stated infinite-rank conclusion in all higher odd
degrees. The construction here uses the displayed arithmetic local ring;
the complete rings considered below use a separate continuous detector.

\section{Visible coefficient rank and all higher odd degrees}
\label{d2ext:dimtwoVisibleHigherRank}

The elliptic contraction detects the $e_3$-component of the point
coefficient. We construct arbitrarily large families whose projections
to this component are $\mathbf Q_5$-linearly independent, and then use
them on a single complete local ring.
We write $\mathcal K(-)=K(-)^\wedge_5[1/5]$ for spectral
$5$-completion followed by inversion of $5$.

For the continuous detector, write $W'=W(\overline{\mathbf F}_5)[\rho]$
for the coefficient ring denoted by $R$ in
Proposition~\ref{d2:prop:weighted}, and put $E_\infty=W'[1/5]$.
The automorphisms $\sigma_g$ fix $W(\overline{\mathbf F}_5)$.
We use the normalization of \eqref{d2:eq:linear-weight}: the bounded
functional has $\ell(\omega)=1$, and the polarization fixes the common
sign. Trace, contraction, and formal replacement are $E_\infty$-linear
on the differential quotients. The same four-term calculation therefore
gives, for every $b\in E_\infty$ with $\sigma_{-1}b=-b$,
\[
 L\operatorname{Tr}_h\left(\sum_{g\in\mathbf F_5^\times}
                  \sigma_g(b)\gamma_g\right)
       =4e_3(b)[\widetilde Q^*\eta].
\]
Here the operators on the scaled model are transported by formal
replacement, as in the proof of Proposition~\ref{d2:prop:weighted}.
The vector $[\widetilde Q^*\eta]$ is nonzero in the ordinary Cohen
quotient by Proposition~\ref{d2:prop:graph}. This formula applies to
coefficients in the finite local fields below and requires only their
$e_3$-projections. It uses the linear operators over the fixed scalar
field, with no passage to a colimit of completed $K$-theory.

For each complete coefficient discrete valuation ring $O$ below and
$n\geq2$, let $r_{n,O}(\beta)$ denote the scalar whose AMMN point
character is $t^{1-n}r_{n,O}(\beta)$, with $|t|=-2$, using the point
comparison in Lemma~\ref{d3:lem:formalhigherTCcup}.

\begin{lemma}
\label{d2ext:dimtwoVisibleHigherCoefficientRank}
Fix $n\geq2$ and $m\geq2$. Put
\[
 V_m=W(\mathbf F_{5^m}),\quad F_m=V_m[1/5],\quad
 O_m=V_m[\rho],\quad E_m=F_m(\rho),\quad \pi=\rho-1,
\]
where $\rho$ is a primitive fifth root of unity.
Let $\sigma_g$ fix $F_m$ and send $\rho$ to $\rho^g$.
Let $\omega_T:\mathbf F_5^\times\to\mathbf Z_5^\times$ be the
Teichm\"uller character, and put
\[
 e_3=\frac14\sum_{g=1}^4\omega_T(g)^{-3}\sigma_g,\qquad
 \tau_3=\sum_{g=1}^4\omega_T(g)^{-3}\rho^g.
\]
There are actual classes
$\beta_j\in K_{2n-1}(O_m)$, indexed by $1\leq j\leq m-1$, each of the form
$(1-\sigma_{-1})\beta_{0,j}$, with the following properties:
their reductions in $K_{2n-1}(\mathbf F_{5^m})$ vanish integrally,
and their AMMN point coordinates $b_j=r_{n,O_m}(\beta_j)\in E_m$
have $\mathbf Q_5$-linearly independent projections $e_3b_j$.
\end{lemma}

\begin{proof}
Define
$S_j(z)=(z\,d/dz)^j(1-z)^{-1}\in\mathbf Z[z,(1-z)^{-1}]$.
The expansion and continuation argument in
the proof of Proposition~\ref{d2:prop:coefficient}, now with
$\binom{n+j-1}{j}$ in place of $j+1$, gives
\[
 F_n(z)\coloneqq \operatorname{Li}_n(z)-5^{-n}\operatorname{Li}_n(z^5)
 =\sum_{j\geq0}(-1)^j\binom{n+j-1}{j}5^j S_j(z^5)
             \sum_{a=1}^4\frac{z^a}{a^{n+j}}.
\]
The series is uniformly convergent on $|z|\leq1$,
$|1-z^5|=1$, and agrees there with Coleman continuation by
\cite[Theorem~2.7(3)]{BesserDeJeuAlgorithm}.
The fifth power acts on the polylogarithm variable.

For a Teichm\"uller root
$\zeta\in\mu_{5^m-1}(V_m)\setminus\{1\}$, the value
$(\zeta\rho)^5=\zeta^5$ is fixed by every $\sigma_g$.
The exact Gauss identity
$e_3(\rho^a)=\frac14\omega_T(a)^3\tau_3$ therefore gives
\begin{equation}
\label{d2ext:dimtwoVisibleHigherDepletion}
\begin{split}
 U_n(\zeta)
 &\coloneqq \frac{4e_3\operatorname{Li}_n(\zeta\rho)}{\tau_3}\\
 &=\sum_{j\geq0}(-1)^j\binom{n+j-1}{j}5^j S_j(\zeta^5)
        \sum_{a=1}^4\frac{\zeta^a\omega_T(a)^3}{a^{n+j}}
       \in V_m.
\end{split}
\end{equation}
The terms with $j\geq1$ belong to $5V_m$.  Thus, writing
$t=\overline\zeta$, its reduction is
\[
 h_{3,n}(t)
   =\frac{P_n(t)}{1-t^5},\qquad
 P_n(t)=\sum_{a=1}^4a^{3-n}t^a\in\mathbf F_5[t],
 \qquad t\in\mathbf F_{5^m}^{\times}\setminus\{1\}.
\]
The polynomial $P_n$ is nonzero, has coefficient one at $t$,
and has degree four.  Hence $h_{3,n}$ is nonconstant and has
rational-map degree at most five.  Every finite fiber has at most
five elements, so
\[
 \#h_{3,n}\bigl(\mathbf F_{5^m}^{\times}\setminus\{1\}\bigr)
 \geq\left\lceil\frac{5^m-2}{5}\right\rceil=5^{m-1}.
\]
Its $\mathbf F_5$-span has dimension at least $m-1$.
Choose $m-1$ independent image values and corresponding
Teichm\"uller roots $\zeta_j$.  Reduction of a normalized
$\mathbf Q_5$-linear relation shows that the $U_n(\zeta_j)$
are $\mathbf Q_5$-independent.  Since $\tau_3\ne0$, the values
$e_3\operatorname{Li}_n(\zeta_j\rho)$ are independent as well.

These are regulators of actual rational number-field classes.
Indeed, all $\zeta_j\rho$ lie in
$L=\mathbf Q(\mu_{5^m-1},\rho)$ and are roots of unity.
At the chosen prime above $5$ they and their complements are
units, since $\overline\zeta_j\ne1$.  The rational symbols
$[\zeta_j\rho]_n$ are cycles, and
\cite[Theorem~1.10(2)]{BesserDeJeu} maps them to a common
nonzero rational multiple of
$L_{\mathrm{mod},n}(\zeta_j\rho)
 =\operatorname{Li}_n(\zeta_j\rho)$.
The point normalization in \eqref{d3:eq:pointnormalization}
identifies the normalized Besser--de Jeu coordinate with the
ordinary syntomic coordinate, so the scalar is common to the
chosen roots. The number-field case of the cited theorem applies
unconditionally. The local completion at the chosen prime is $E_m$.

Localization at that coefficient prime lifts these rational
classes to its integer ring, since
$K_{2n-2}(\mathbf F_{5^m})=0$.
Clear one common integer denominator and then form the actual
differences
$\beta_j=(1-\sigma_{-1})\beta_{0,j}$, mapping them to $K_{2n-1}(O_m)$.
Their residue reductions are zero because $\sigma_{-1}$ fixes
the residue field.  Applying $e_3$ to their syntomic regulators
multiplies each of the preceding values by the same nonzero
rational scalar, including the factor two from the difference.

The ramified point comparison
\cite[Example~2.2.4 and Propositions~2.2.7, 2.2.9, 2.3.4]{HuberKings}
is Galois-equivariant and identifies syntomic and \'etale point
cohomology in every twist $n\geq2$.  Thus the $e_3$-projections
of the \'etale Chern values are $\mathbf Q_5$-independent.
Lemma~\ref{d3:lem:cherncompletionfactor} gives a natural
Galois-equivariant $\mathbf Q_5$-linear map
\[
 \pi_{2n-1}\bigl(K(O_m)^\wedge_5[1/5]\bigr)
       \longrightarrow H^1(E_m,\mathbf Q_5(n))
\]
through which the actual Chern class factors. The finite Chern
maps are applied over $E_m$, where $5$ is invertible, after the
canonical projections to $K_{2n-1}(O_m;\mathbf Z/5^\nu)$ and
restriction to the field. Compatibility with reduction gives the
map to the inverse limit of $H^1$. The finite $H^0$ groups are
Mittag--Leffler, identifying this target with continuous cohomology.
This construction uses the canonical source projections and
allows a nonzero source $\lim^1$.

The point case of the AMMN comparison identifies the source with
$t^{1-n}E_m$, naturally under $\sigma_g$.  A relation among the
$e_3b_j$ would therefore induce the same relation among the
independent projected Chern values.  This proves the claim
without identifying the AMMN and polylogarithm coordinates.
\end{proof}

\begin{remark}
The rational-map degree bound cannot uniformly be replaced by four.
When $n\equiv3\pmod4$ one has
\[
 h_{3,n}(t)=(t-1)^{-5}-(t-1)^{-1}.
\]
Its image lies in the trace-zero hyperplane of
$\mathbf F_{5^m}/\mathbf F_5$.  The lower bound $m-1$ accommodates
this actual loss of one dimension.
\end{remark}

Put
\[
 k=\overline{\mathbf F}_5,\qquad W=W(k),\qquad
 V=\widehat{W[T]_{(5)}}^{\,(5)},\qquad
 R=V[[x,y]]/(\Phi_5(1+y)+x^4+Txy^3).
\]

\begin{theorem}
\label{d2ext:dimtwoCompleteAllOddInfiniteRank}
The ring $R$ is a complete regular local domain of dimension two,
with maximal ideal $(x,y)$ and residue field $k(T)$.
For every $n\geq2$,
\[
 \dim_{\mathbf Q}\ker\left(
 K_{2n-1}(R)_{\mathbf Q}\longrightarrow
 K_{2n-1}(\operatorname{Frac}R)_{\mathbf Q}\right)=\infty.
\]
For each fixed $n,m\geq2$, the kernel contains $m-1$ independent
classes represented by actual integral classes that vanish on
$R[1/x]$ integrally.  The ring $R$ is independent of $n,m$.
\end{theorem}

\begin{proof}
The regularity argument is the same linear-term calculation as in
Section~\ref{d2:sec:preliminaries}.
Set $W'=W[\rho]$ and $E_\infty=W'[1/5]$.
For fixed $n,m$, take the actual coefficients of
Lemma~\ref{d2ext:dimtwoVisibleHigherCoefficientRank}.
Naturality of the AMMN point calculation sends their coordinates
under $E_m\hookrightarrow E_\infty$; hence their visible
projections remain $\mathbf Q_5$-independent.

Use the same finite parameter cover and scaled formal test ring
$B$ as in \eqref{d2:eq:test-algebra}, now over $V$.
It is the completion of an essentially smooth $W'$-algebra
of relative dimension two.  The spectrum comparison underlying
Lemma~\ref{d3:lem:essentiallysmoothAMMN} gives, for each fixed $n$,
the relative character to $HC_{2n-2}^{\mathrm{cts}}(B/W';\mathbf Q_5)$.
In its finite mixed-HKR model the top-form summand is
\begin{equation}
\label{d2ext:dimtwoHigherTopColumn}
 t^{2-n}\left(
 \Omega^2_{B/W'}[1/5]/d\Omega^1_{B/W'}[1/5]\right).
\end{equation}
Indeed its cycles in this summand are all two-forms, and its
boundaries are precisely $t^{2-n}d\Omega^1$; form degrees above
two vanish.

The zero residue reductions of the coefficients provide actual
relative lifts at the point. Pull them to $B$ and multiply by
$\eta_g=1-[J_g]$.
With the common scalar $\lambda_{W'}\in E_\infty^\times$ of
Proposition~\ref{d2:prop:cyclic}, the negative-cyclic module product is
\[
 (t^{1-n}\sigma_g(b_j))\,(\lambda_{W'}t\gamma_g)
       =\lambda_{W'}t^{2-n}\sigma_g(b_j)\gamma_g.
\]
Rank zero removes the scalar term, higher Chern forms vanish in
dimension two, and the cyclic-shuffle terms vanish because the
point coefficient is a scalar relative to $W'$.
The scalar-bar calculation applies in every cyclic column, including
the negative columns that occur for $n>2$.
The scalar occurs to the first power because it belongs to the
first Chern component.

Trace and contraction in Proposition~\ref{d2:prop:weighted} are
$E_\infty$-linear. They send the coefficient of
\eqref{d2ext:dimtwoHigherTopColumn} to
\[
 4\lambda_{W'}e_3(b_j)[\widetilde Q^*\eta].
\]
The fixed vector $[\widetilde Q^*\eta]$ is nonzero over $E_\infty$
after the generic Cohen completion, by
Lemma~\ref{d2:lem:scalar} and Proposition~\ref{d2:prop:graph}.
Multiplication by the common nonzero $\lambda_{W'}$ is an
injective $\mathbf Q_5$-linear map on this coefficient field.
Thus these character images are $\mathbf Q_5$-independent,
without any rationality or Galois-invariance assumption on
$\lambda_{W'}$.

Write $C=B/\pi$.  As before, $C$ is a filtered colimit of smooth
affine surfaces over $k$.  By
\cite[Theorem~8.4]{GeisserLevineCharP},
$\pi_q\mathcal K(C)=0$ for every $q\geq3$.
The relative and regular-divisor localization sequences therefore
give, in the present degree,
\[
 \pi_{2n-1}\mathcal K(B,(\pi))
 \hookrightarrow\pi_{2n-1}\mathcal K(B)
 \hookrightarrow\pi_{2n-1}\mathcal K(B[1/\pi]).
\]
These use the vanishings in degrees $2n$ and $2n-1$, respectively.
These are vanishings of the rationalized completed spectrum.

The source divisor $x=0$ has ring $V[\rho]$.
Pull each $\beta_j$ to that divisor and push forward to define
$c_j\in K_{2n-1}(R)$.  Localization gives $c_j|_{R[1/x]}=0$
integrally.  Finite flat coefficient extension $V\to V[\rho]$,
followed by inversion of $5$, splits the divisor into the four
root components with coefficients $\sigma_g(\beta_j)$.
Pullback of their line-bundle classes under the scaled map gives
\[
 c_j|_{B[1/\pi]}
   =\sum_g\sigma_g(\beta_j)(1-[J_g])\big|_{B[1/\pi]}.
\]
This is exactly the actual product detected above; the scaled map
itself is not required to be flat.
Consequently the $c_j$ are rationally independent.
Since $m$ is arbitrary while $R$ is fixed, the kernel has infinite
dimension for this $n$.  Applying the argument separately to each
$n\geq2$ proves the theorem.
\end{proof}

\section{One complete two-dimensional ring in every degree at least three}
\label{d2ext:dimtwoAllDegreesFixedRing}

The additional parameter below is a unit of the source ring.
Its logarithmic differential raises the degree of the detected
class by one.  Detecting the even products also proves independence
of their odd factors, including degree three.

\begin{lemma}
\label{d2ext:dimtwoFullCohenLaurentResidue}
Let $W'$ be the coefficient discrete valuation ring, with
uniformizer $\pi$, used in the quartic construction.  Let $D_0$
be a smooth affine $W'$-curve whose special fiber is integral,
and put
\[
 A_T=\widehat{(D_0)_{(\pi)}}^{\,(\pi)},\qquad
 A_{TS}=\widehat{(D_0[S])_{(\pi)}}^{\,(\pi)}.
\]
Here localization at $(\pi)$ inverts every element with nonzero
reduction.  Let $B_T$ be the $\pi$-adic completion of a smooth
affine relative curve over $A_T$, defined by a model over $D_0$,
and let $B_{TS}$ be the corresponding completed model over
$A_{TS}$.  Suppose their continuous differentials over $W'$
are finite projective of ranks two and three, respectively.
There is a $W'$-linear operator
\[
 \operatorname{Res}_S:
 \Omega^{r,\mathrm{cts}}_{B_{TS}/W'}
       \longrightarrow \Omega^{r-1,\mathrm{cts}}_{B_T/W'}
\]
which is integral and $\pi$-adically continuous and satisfies
\[
 \operatorname{Res}_S(d\alpha)=-d\operatorname{Res}_S(\alpha),
 \qquad
 \operatorname{Res}_S(\theta\wedge d\log S)=\theta
       \quad\text{for every }\theta\in\Omega^2_{B_T/W'}.
\]
In particular it sends actual exact three-forms to actual exact
two-forms. The induced map
\[
 \frac{\Omega^2_{B_T/W'}[1/5]}
      {d\Omega^1_{B_T/W'}[1/5]}
 \longrightarrow
 \frac{\Omega^3_{B_{TS}/W'}[1/5]}
      {d\Omega^2_{B_{TS}/W'}[1/5]},
 \qquad [\theta]\longmapsto[\theta\wedge d\log S],
\]
is injective; these are ordinary continuous quotients.
\end{lemma}

\begin{proof}
For every $a\geq1$ there is a compatible ring map
\[
 B_{TS}/\pi^a
       \longrightarrow (B_T/\pi^a)((S)).
\]
Indeed, if $g\in D_0[S]$ has nonzero reduction, its image modulo
$\pi$ is a nonzero polynomial over the field
$\operatorname{Frac}(D_0/\pi)$.  It is therefore a unit in that
field's Laurent series ring.  The ideal generated by $\pi$ in
$(B_T/\pi^a)((S))$ is nilpotent, so this inverse lifts.  This
argument applies to every primitive denominator, including
denominators whose literal lowest coefficient is not a unit.
The lifts of the inverses are unique and commute with reduction
in $a$.  The maps extend to the completed algebra because the
inverse-limit target is complete.  Thus the relevant target is
\[
 \mathcal L=\varprojlim_a(B_T/\pi^a)((S)).
\]
The negative Laurent bound may depend on $a$, which is why this
inverse-limit target is used.

Use Laurent series of continuous base forms, with differential
defined coefficientwise and with the usual $S$-derivative.
These operations are compatible modulo $\pi^a$ and preserve
the $\pi$-adic filtration.  The universal continuous derivation
of $B_{TS}$ therefore gives a map into this Laurent differential
complex.  Finite projectivity of the source form modules
justifies extension to their completions and exterior powers.

Write a form of degree $r$ in the target as
$\alpha(S)+\beta(S)\wedge dS$, with base-form degrees $r$
and $r-1$.  Define
\[
 \operatorname{Res}_{S,r}
       \bigl(\alpha(S)+\beta(S)\wedge dS\bigr)
      =(-1)^{r-1}[S^{-1}]\beta(S).
\]
Coefficient extraction is compatible with reduction and is
integral.  The coefficient of $S^{-1}$ in the derivative of a
Laurent series is zero.  It follows directly that
$\operatorname{Res}_{S,r+1}d=-d\operatorname{Res}_{S,r}$.
For $r=3$ its sign is positive, which gives the second displayed
identity.  Inverting $5$ preserves these identities.  All
quotients in the conclusion use the ordinary image of the
differential.
\end{proof}

Put
\[
 k=\overline{\mathbf F}_5,\qquad W_0=W(k),\qquad
 V_{TS}=\widehat{W_0[T,S]_{(5)}}^{\,(5)},
\]
and
\[
 R_{TS}=V_{TS}[[x,y]]/
       \bigl(\Phi_5(1+y)+x^4+Txy^3\bigr).
\]

\begin{theorem}
\label{d2ext:dimtwoCompleteAllDegreesInfiniteRank}
The ring $R_{TS}$ is a complete regular local domain of dimension
two, with maximal ideal $(x,y)$ and residue field $k(T,S)$.
For every integer $j\geq3$,
\[
 \dim_{\mathbf Q}\ker\!\left(
 K_j(R_{TS})_{\mathbf Q}\longrightarrow
 K_j(\operatorname{Frac}R_{TS})_{\mathbf Q}\right)=\infty.
\]
For every $n,m\geq2$, the odd degree $2n-1$ and the even degree
$2n$ kernels contain $m-1$ independent classes represented by
actual integral classes which vanish on $R_{TS}[1/x]$ integrally.
The ring $R_{TS}$ is independent of $n$ and $m$.
\end{theorem}

\begin{proof}
The ring $W_0[T,S]_{(5)}$ is a discrete valuation ring with
uniformizer $5$ and residue field $k(T,S)$.  Its completion
$V_{TS}$ has the same properties.  The defining equation of
$R_{TS}$ is congruent to $5$ modulo $(5,x,y)^2$ in the regular
local ambient ring of dimension three.  Its quotient is
therefore regular of dimension two, and is complete with
maximal ideal $(x,y)$.  In particular it is a domain.
The element $S$ is a unit of $R_{TS}$.

Set $W'=W_0[\rho]$, $\pi=\rho-1$, and
$E_\infty=W'[1/5]$.  Fix $n,m\geq2$, and take the $m-1$
actual coefficients $\beta_j$ from
Lemma~\ref{d2ext:dimtwoVisibleHigherCoefficientRank}.
Map them to $K_{2n-1}(W')$ and write $b_j\in E_\infty$
for their natural AMMN point coordinates.  Their visible
projections $e_3b_j$ remain $\mathbf Q_5$-independent.
Their reductions vanish integrally, since each coefficient
is an actual difference $(1-\sigma_{-1})\beta_{0,j}$.

Choose the same finite parameter cover, independent of $S$,
as in the one-parameter quartic detector.  In the notation of
Lemma~\ref{d2ext:dimtwoFullCohenLaurentResidue}, it gives the old
coefficient ring $A_T$ and the new one $A_{TS}$.  The latter
receives $V_{TS}[\rho]$: every primitive polynomial in $T,S$
has nonzero image in $k(U)(S)$, and hence is a unit.  Completion
extends this map.  Put
\[
 B_{TS}=A_{TS}\langle X,Y\rangle/
 \left(X^4+\frac{\Phi_5(1+\pi Y)}{\pi^4}+TXY^3\right),
 \qquad C_{TS}=B_{TS}/\pi.
\]
The old algebra $B_T$ is given by the same formula over $A_T$.
They are the completions of essentially smooth $W'$-algebras
of relative dimensions three and two, respectively.
Smoothness holds near the entire special fiber; localizing a
model at an element congruent to one modulo $\pi$ leaves the
completion unchanged.  In particular $B_{TS}$ and $C_{TS}$
are regular and $\pi$ is a non-zero-divisor.

For $g\in\mathbf F_5^\times$, let
\[
 a_g=\frac{\rho^g-1}{\pi},\qquad
 J_g=(X,Y-a_g),\qquad
 \eta_g=1-[J_g],\qquad \gamma_g=-c_1(J_g).
\]
These invertible ideals are pulled from $B_T$.  We can therefore
represent their Chern classes by two-forms pulled from the old
continuous complex, independently of $S$.

The actual residue sequence at the point supplies actual
relative lifts of all $\beta_j$; here
$K_{2n}(k)=0$ by Quillen's finite-field computation and filtered
colimits.  Pull these lifts to $B_{TS}$ and form their relative
products
\[
 \xi_j=\sum_g\sigma_g(\beta_j)\eta_g[S]
       \quad\text{in }K_{2n}(B_{TS},(\pi)),
\]
where the coefficient symbols in this formula denote the
specified relative lifts.

Apply the essentially smooth relative fiber comparison of
Lemma~\ref{d3:lem:essentiallysmoothAMMN} in each fixed degree.
Its $K(B_{TS})$-module compatibility gives the negative-cyclic
action on the ordinary cyclic target.  In relative dimension
three the normalized finite mixed-HKR map uses only the
denominators $2$ and $6$, both units at $5$.
Let $\lambda_{W'}\in E_\infty^\times$ be the common scalar of
Proposition~\ref{d2:prop:cyclic}.
The line bundles are pulled from the rank-two algebra $B_T$,
so naturality pulls back their negative-cyclic characters
$\lambda_{W'}t\gamma_g$.
The point coefficient is $t^{1-n}\sigma_g(b_j)$, and the
unit character starts with $d\log S$.
The top character component is
\[
 \chi_{B_{TS}}(\xi_j)_{\mathrm{top}}
    =\lambda_{W'}t^{2-n}
       \sum_g\sigma_g(b_j)\gamma_g\wedge d\log S
       \quad\text{in }HC^{\mathrm{cts}}_{2n-1}(B_{TS}/W';
                                                    \mathbf Q_5).
\]
Here $|t|=-2$.  This formula also holds when the initial cyclic
column is negative.  To check it at chain level, the rank-zero
line has underlying Hochschild degrees at least two; the
lowest unit term has degree one.  Every higher line or unit
term, or cyclic-shuffle correction in their product, has final
Hochschild degree at least five, and is killed by the
dimension-three HKR map.  The scalar point cup has no
cyclic-shuffle correction because its scalar lies in the
relative base.  Thus precisely the displayed ordinary
shuffle survives.

The top-form summand is
\[
 t^{2-n}\bigl(
 \Omega^3_{B_{TS}/W'}[1/5]/
                  d\Omega^2_{B_{TS}/W'}[1/5]\bigr).
\]
All its three-forms are cycles, and its boundaries are precisely
the indicated ordinary exact forms.  Apply
Lemma~\ref{d2ext:dimtwoFullCohenLaurentResidue} to their coefficients.
It recovers
$\lambda_{W'}\sum_g\sigma_g(b_j)\gamma_g$ on $B_T$.
The weighted trace and elliptic contraction of
Proposition~\ref{d2:prop:weighted} are $E_\infty$-linear and send
these classes to
\[
 4\lambda_{W'}e_3(b_j)[\widetilde Q^*\eta].
\]
The fixed base vector $[\widetilde Q^*\eta]$ survives the old generic
Cohen completion by Lemma~\ref{d2:lem:scalar}, and remains
nonzero after the finite constant extension adjoining $\rho$ by
Proposition~\ref{d2:prop:graph}.
Multiplication by the common nonzero $\lambda_{W'}$ preserves
$\mathbf Q_5$-linear independence.
Consequently the characters of the $\xi_j$ are
$\mathbf Q_5$-independent.

The special fiber $C_{TS}$ is a filtered colimit of smooth
affine threefolds over the perfect field $k$.  Geisser--Levine
\cite[Theorem~8.4]{GeisserLevineCharP} gives
\[
 K_q(C_{TS};\mathbf Z/5^\nu)=0
      \quad(q>3,\ \nu\geq1).
\]
The Milnor sequence for spectral completion, using degrees
$q$ and $q+1$, therefore gives
$\pi_q\mathcal K(C_{TS})=0$ for $q\geq4$.
For the even degree under consideration the relative and
regular-divisor localization sequences give
\[
 \pi_{2n}\mathcal K(B_{TS},(\pi))
   \hookrightarrow\pi_{2n}\mathcal K(B_{TS})
   \hookrightarrow\pi_{2n}\mathcal K(B_{TS}[1/\pi]).
\]
The first injection uses the vanishing in degree $2n+1$,
and the second that in degree $2n$; both are valid for $n\geq2$.
Here $\mathcal K=K^\wedge_5[1/5]$ denotes the rationalized
spectrally completed $K$-theory.

It remains to match these detected even classes with actual
classes of the fixed source.  Its regular Cartier divisor
$i:x=0\hookrightarrow\operatorname{Spec}R_{TS}$ has ring
$V_{TS}[\rho]$.  Pull the coefficients to this divisor and set
\[
 c_{n,j}=i_*(\beta_j)\in K_{2n-1}(R_{TS}),\qquad
 a_{n,j}=c_{n,j}[S]\in K_{2n}(R_{TS}).
\]
Both classes vanish on $R_{TS}[1/x]$ integrally.
Make the finite flat coefficient extension
$V_{TS}\to V_{TS}[\rho]$, and then invert $5$.
The divisor $x=0$ splits into its four root components with
coefficients $\sigma_g(\beta_j)$.
Their length-one Cartier resolutions and the projection
formula express the resulting odd class as
$\sum_g\sigma_g(\beta_j)(1-[I_g])$, where
$I_g=(x,y-(\rho^g-1))$.
Apply the continuous scaling
\[
 x\longmapsto\pi X,\qquad y\longmapsto\pi Y,\qquad
 R_{TS}\longrightarrow B_{TS}.
\]
It exists because the defining equation maps to $\pi^4$ times
the test equation and the maximal ideal maps into $(\pi)$.
After inverting $\pi$, pullback of the invertible ideals $I_g$
gives $J_g$. The resulting actual class is
\[
 a_{n,j}|_{B_{TS}[1/\pi]}
   =\sum_g\sigma_g(\beta_j)\eta_g[S]\big|_{B_{TS}[1/\pi]}
       \quad\text{in actual }K_{2n}.
\]
Only the initial coefficient extension is asserted to be flat;
the subsequent substitution pulls back line bundles.
This equality and the two completed injections prove that
the actual classes $a_{n,j}$ are rationally independent.

Finally, a rational relation among the actual odd classes
$c_{n,j}$ would, after multiplication by the same actual unit
class $[S]$ in the same ring, give a relation among the
$a_{n,j}$.  Thus the odd classes are independent as well.
This argument includes $n=2$ and requires no vanishing of
$\pi_3\mathcal K(C_{TS})$.
Letting $m$ grow proves infinite rank separately in degrees
$2n-1$ and $2n$.  These pairs cover every $j\geq3$.
\end{proof}

\begin{remark}
No general injectivity of multiplication by $[S]$ is used:
it is proved only on the displayed finite spans by detecting
their products.  In particular the odd degree-three conclusion
does not use a degree-three generic-localization injection for
the new special fiber.  The residue field in this fixed-ring
statement is $\overline{\mathbf F}_5(T,S)$.
\end{remark}

